\begin{tuluyen}
Trong mỗi trường hợp sau, chỉ rõ yếu tố nào ảnh hưởng đến tốc độ phản ứng?
\begin{enumerate}
\item Thanh củi được chẻ nhỏ hơn thì sẽ cháy nhanh hơn.
\item Khi bảo quản thực phẩm trong tủ lạnh thì thực phẩm tươi lâu hơn khi để trong không khí.
\end{enumerate}
\source{Trích nguyên văn Câu 2.1, đề tuyển sinh chuyên Hóa Đắk Lắk 2026.}
\dapso{(a) Diện tích bề mặt; (b) nhiệt độ.}
\begin{loigiaiBT}
Khi chẻ nhỏ củi, tổng diện tích tiếp xúc với oxygen tăng nên phản ứng cháy diễn ra nhanh hơn.

Khi bảo quản thực phẩm trong tủ lạnh, nhiệt độ thấp làm chậm các phản ứng hóa học và quá trình sinh hóa gây hư hỏng.

\textbf{Vậy,} hai yếu tố tương ứng là diện tích bề mặt và nhiệt độ.
\end{loigiaiBT}
\end{tuluyen}
